\documentclass[11pt]{article}

\usepackage{iftex}
\usepackage[spanish]{babel}
\usepackage{amsmath,amssymb,amsthm,mathtools}
\usepackage{geometry}
\usepackage{enumitem}

\geometry{margin=1in}
\setlist[enumerate]{label=\textup{(\roman*)}}

\ifPDFTeX
  \usepackage[utf8]{inputenc}
  \usepackage[T1]{fontenc}
\else
  \usepackage{fontspec}
  \usepackage{unicode-math}
  \defaultfontfeatures{Ligatures=TeX}
  \setmainfont{texgyrepagella-regular.otf}[
    BoldFont=texgyrepagella-bold.otf,
    ItalicFont=texgyrepagella-italic.otf,
    BoldItalicFont=texgyrepagella-bolditalic.otf
  ]
  \setsansfont{texgyreheros-regular.otf}[
    BoldFont=texgyreheros-bold.otf,
    ItalicFont=texgyreheros-italic.otf,
    BoldItalicFont=texgyreheros-bolditalic.otf
  ]
  \setmonofont{DejaVuSansMono.ttf}[
    BoldFont=DejaVuSansMono-Bold.ttf,
    Scale=MatchLowercase
  ]
  \setmathfont{texgyrepagella-math.otf}
\fi

\newtheorem{theorem}{Teorema}
\newtheorem{lemma}[theorem]{Lema}
\newtheorem{proposition}[theorem]{Proposici\'on}
\newtheorem{corollary}[theorem]{Corolario}
\theoremstyle{definition}
\newtheorem{definition}[theorem]{Definici\'on}
\newtheorem{example}[theorem]{Ejemplo}
\newtheorem{exercise}[theorem]{Ejercicio}
\theoremstyle{remark}
\newtheorem{remark}[theorem]{Observaci\'on}

\newcommand{\R}{\mathbb{R}}
\newcommand{\norm}[1]{\left\lVert #1 \right\rVert}
\newcommand{\dd}{\mathrm{d}}

\title{1er Teorico Geometria Superior}
\author{EDISON ALBERTO FERNANDEZ CULMA}
\date{}

\begin{document}
\maketitle

\section*{Diferenciabilidad en espacios normados}

\begin{definition}
Sean $V$ y $W$ espacios vectoriales reales de dimensi\'on finita, normados. Sea $I \subseteq V$ un abierto y $p \in I$. Una funci\'on $F \colon I \to W$ se dice \emph{diferenciable en $p$} si existe una transformaci\'on lineal $T \colon V \to W$ tal que
\[
\lim_{v \to 0} \frac{\norm{F(p+v)-F(p)-Tv}}{\norm{v}} = 0.
\]
La norma del denominador es la de $V$ y la del numerador es la de $W$.
\end{definition}

\paragraph{Pregunta.}
\`{Que} representa $Tv$?

\begin{proposition}
Sea $F \colon I \subseteq V \to W$ diferenciable en $p$ y sea $T \colon V \to W$ una transformaci\'on lineal tal que
\[
\lim_{v \to 0} \frac{\norm{F(p+v)-F(p)-Tv}}{\norm{v}} = 0.
\]
Entonces, para todo $v \in V$,
\[
Tv = \lim_{t \to 0} \frac{F(p+tv)-F(p)}{t}.
\]
\end{proposition}

\begin{proof}
Si $v=0$, la igualdad es inmediata porque $T0=0$. Supongamos $v \neq 0$. Como $I$ es abierto y $p \in I$, existe $\varepsilon>0$ tal que $p+tv \in I$ para todo $t \in (-\varepsilon,\varepsilon)$. Entonces
\[
0
= \lim_{t \to 0}
\frac{\norm{F(p+tv)-F(p)-T(tv)}}{\norm{tv}}
= \lim_{t \to 0}
\frac{1}{\norm{v}}
\norm{\frac{F(p+tv)-F(p)}{t}-Tv}.
\]
Multiplicando por $\norm{v}$, se obtiene
\[
\lim_{t \to 0}\norm{\frac{F(p+tv)-F(p)}{t}-Tv}=0,
\]
de donde
\[
Tv = \lim_{t \to 0} \frac{F(p+tv)-F(p)}{t}.
\qedhere
\]
\end{proof}

\begin{corollary}
Si $F$ es diferenciable en $p$, existe una \'unica transformaci\'on lineal que satisface la definici\'on anterior. Se la llama \emph{diferencial de $F$ en $p$} y se la denota por
\[
(\dd F)_p \colon V \to W.
\]
\end{corollary}

\begin{proof}
La proposici\'on anterior muestra que cualquier transformaci\'on lineal que satisfaga la definici\'on debe verificar
\[
Tv = \lim_{t \to 0} \frac{F(p+tv)-F(p)}{t}
\qquad \text{para todo } v \in V,
\]
y por lo tanto queda completamente determinada por $F$ y por $p$.
\end{proof}

De la definici\'on de diferencial se obtiene la descomposici\'on
\[
F(p+v)=F(p)+(\dd F)_p(v)+R(v),
\qquad
\lim_{v\to 0}\frac{\norm{R(v)}}{\norm{v}}=0.
\]
El t\'ermino $R(v)$ es el error al aproximar $F(p+v)$ por la expresi\'on af\'in
\[
F(p)+(\dd F)_p(v).
\]

\begin{definition}
Si $F \colon I \subseteq V \to W$ es diferenciable en $p$, la funci\'on
\[
A_p(x)=F(p)+(\dd F)_p(x-p)
\]
se llama \emph{aproximaci\'on af\'in de $F$ en $p$}.
\end{definition}

\begin{lemma}
Toda transformaci\'on lineal $T \colon V \to W$ entre espacios vectoriales reales normados de dimensi\'on finita es acotada. Es decir, existe $\mu \ge 0$ tal que
\[
\norm{Tv}\le \mu \norm{v}
\qquad \text{para todo } v \in V.
\]
\end{lemma}

\begin{proof}
Sea
\[
S=\{v\in V:\norm{v}=1\}.
\]
La funci\'on $G \colon S \to \R$, $G(v)=\norm{Tv}$, es continua. Como $S$ es cerrado y acotado en un espacio de dimensi\'on finita, es compacto; por ende $G(S)$ alcanza un m\'aximo, digamos $\mu\ge 0$. Luego
\[
\norm{Tv}\le \mu
\qquad \text{para todo } v\in S.
\]
Si $v\neq 0$, entonces $v/\norm{v}\in S$, de modo que
\[
\norm{Tv}
\;=\;
\norm{v}\,\norm{T\!\left(\frac{v}{\norm{v}}\right)}
\;\le\;
\mu \norm{v}.
\]
El caso $v=0$ es trivial.
\end{proof}

\begin{theorem}
Si $F \colon I \subseteq V \to W$ es diferenciable en $p$, entonces $F$ es continua en $p$.
\end{theorem}

\begin{proof}
Sea $T=(\dd F)_p$. Entonces
\[
F(p+v)-F(p)=\bigl(F(p+v)-F(p)-Tv\bigr)+Tv.
\]
Tomando normas y usando el lema anterior,
\[
\norm{F(p+v)-F(p)}
\le
\norm{F(p+v)-F(p)-Tv}+\norm{Tv}
\le
\norm{v}\,\frac{\norm{F(p+v)-F(p)-Tv}}{\norm{v}}+\mu \norm{v}.
\]
Cuando $v\to 0$, ambos sumandos tienden a $0$, luego
\[
\lim_{v\to 0} F(p+v)=F(p).
\]
\end{proof}

\begin{exercise}
Rehacer la demostraci\'on del lema anterior usando continuidad de transformaciones lineales, o bien demostrar que toda transformaci\'on lineal entre espacios normados de dimensi\'on finita es continua.
\end{exercise}

\section*{Diferenciabilidad en conjuntos y regla de la cadena}

\begin{definition}
Sea $F \colon I \subseteq V \to W$ con $I$ abierto. Se dice que $F$ es \emph{diferenciable} si es diferenciable en cada punto de $I$.
\end{definition}

\begin{definition}
Sea $A\subseteq V$ arbitrario y sea $F \colon A \to W$ continua. Se dice que $F$ es \emph{diferenciable en $p\in A$} si existe un abierto $I_p\subseteq V$ con $p\in I_p$ y una funci\'on diferenciable
\[
F_p \colon I_p \to W
\]
tal que $F_p$ extiende a $F$ en $A\cap I_p$, es decir,
\[
F_p|_{A\cap I_p}=F|_{A\cap I_p}.
\]
\end{definition}

\begin{remark}
M\'as adelante se puede probar una caracterizaci\'on equivalente: $F \colon A \to W$ es diferenciable si y solo si existe un abierto $I$ que contiene a $A$ y una funci\'on diferenciable $\widetilde{F}\colon I\to W$ tal que $\widetilde{F}|_A=F$.
\end{remark}

\begin{theorem}[Regla de la cadena]
Sean $V$, $W$ y $X$ espacios vectoriales reales normados de dimensi\'on finita. Sean $I\subseteq V$ y $\widetilde{I}\subseteq W$ abiertos, y sean
\[
F \colon I \to \widetilde{I},
\qquad
G \colon \widetilde{I} \to X.
\]
Si $F$ es diferenciable en $p\in I$ y $G$ es diferenciable en $F(p)\in \widetilde{I}$, entonces $G\circ F$ es diferenciable en $p$ y
\[
\dd(G\circ F)_p
=
\dd G_{F(p)} \circ \dd F_p.
\]
\end{theorem}

\begin{proof}
Escribamos
\[
q=F(p),\qquad T=\dd F_p,\qquad Q=\dd G_q.
\]
Queremos probar que
\[
\lim_{v\to 0}
\frac{\norm{G(F(p+v))-G(F(p))-Q(Tv)}}{\norm{v}}=0.
\]
Definimos
\[
w = F(p+v)-F(p).
\]
Entonces $w\to 0$ cuando $v\to 0$ porque $F$ es continua en $p$. Adem\'as,
\[
G(F(p+v))-G(F(p))-Q(Tv)
=
\bigl(G(q+w)-G(q)-Qw\bigr)+Q(w-Tv).
\]
Por el lema de acotaci\'on aplicado a $Q$, existe $\mu\ge 0$ tal que
\[
\norm{Q(w-Tv)}\le \mu \norm{w-Tv}.
\]
Por lo tanto,
\[
\frac{\norm{G(F(p+v))-G(F(p))-Q(Tv)}}{\norm{v}}
\le
\frac{\norm{G(q+w)-G(q)-Qw}}{\norm{v}}
\;+\;
\mu \frac{\norm{w-Tv}}{\norm{v}}.
\]
Si $w=0$, el primer sumando es cero. Si $w\neq 0$, escribimos
\[
\frac{\norm{G(q+w)-G(q)-Qw}}{\norm{v}}
=
\frac{\norm{G(q+w)-G(q)-Qw}}{\norm{w}}
\cdot
\frac{\norm{w}}{\norm{v}}.
\]
Adem\'as,
\[
\frac{\norm{w}}{\norm{v}}
\le
\frac{\norm{F(p+v)-F(p)-Tv}}{\norm{v}}
\;+\;
\frac{\norm{Tv}}{\norm{v}}
\le
\frac{\norm{F(p+v)-F(p)-Tv}}{\norm{v}}+\lambda
\]
para alguna constante $\lambda\ge 0$, nuevamente por acotaci\'on de $T$. El primer factor tiende a $0$ porque $G$ es diferenciable en $q$, y el segundo est\'a acotado cerca de $v=0$. Por otro lado,
\[
\mu \frac{\norm{w-Tv}}{\norm{v}}
=
\mu \frac{\norm{F(p+v)-F(p)-Tv}}{\norm{v}}
\longrightarrow 0
\qquad (v\to 0)
\]
porque $F$ es diferenciable en $p$. Concluimos que el cociente completo tiende a $0$.
\end{proof}

\section*{Derivadas parciales}

Nos enfocamos ahora en el caso $V=\R^n$ y $W=\R^m$ con sus normas usuales. Sea
\[
F \colon U \subseteq \R^n \to \R^m
\]
diferenciable en $p\in U$. En $\R^n$ consideramos la base can\'onica
\[
\{e_1,\dots,e_n\},
\qquad
e_i=(0,\dots,0,1,0,\dots,0).
\]

La matriz de $(\dd F)_p$ respecto de las bases can\'onicas de $\R^n$ y $\R^m$ se llama \emph{matriz jacobiana} de $F$ en $p$ y se denota por
\[
JF(p)
=
\bigl((\dd F)_p e_1 \ \cdots \ (\dd F)_p e_n\bigr).
\]
Cada columna es el vector $(\dd F)_p e_i$, y por la proposici\'on inicial,
\[
(\dd F)_p e_i
=
\lim_{t\to 0}\frac{F(p+t e_i)-F(p)}{t}.
\]

Si escribimos
\[
(\dd F)_p e_i = \sum_{j=1}^m t_j e_j,
\]
entonces las coordenadas $t_j$ satisfacen
\[
t_j
=
\left\langle (\dd F)_p e_i,e_j\right\rangle
=
\lim_{t\to 0}
\left\langle \frac{F(p+t e_i)-F(p)}{t},e_j\right\rangle.
\]
Si $F=(F_1,\dots,F_m)$, esto equivale a
\[
t_j
=
\lim_{t\to 0}\frac{F_j(p+t e_i)-F_j(p)}{t}.
\]

\begin{definition}[Derivada parcial escalar]
Dada $f \colon U \subseteq \R^n \to \R$, con $U$ abierto, la \emph{$i$-\'esima derivada parcial} de $f$ en $p\in U$ se define por
\[
\frac{\partial f}{\partial x_i}(p)
:=
\lim_{t\to 0}\frac{f(p+t e_i)-f(p)}{t},
\]
cuando este l\'imite existe.
\end{definition}

\begin{definition}[Derivada parcial vectorial]
Si
\[
F=(F_1,\dots,F_m)\colon U \subseteq \R^n \to \R^m,
\]
definimos
\[
\frac{\partial F}{\partial x_i}(p)
:=
\begin{pmatrix}
\dfrac{\partial F_1}{\partial x_i}(p)\\[0.6em]
\vdots\\[0.2em]
\dfrac{\partial F_m}{\partial x_i}(p)
\end{pmatrix},
\]
cuando cada una de esas derivadas parciales existe.
\end{definition}

\begin{remark}
Si $F$ es diferenciable en $p$, entonces la matriz jacobiana est\'a dada por
\[
JF(p)=\bigl(a_{jk}\bigr)
=
\left(\frac{\partial F_j}{\partial x_k}(p)\right),
\]
es decir, la $i$-\'esima columna de $JF(p)$ es el vector
\[
\frac{\partial F}{\partial x_i}(p).
\]
\end{remark}

\begin{definition}[Funci\'on de clase $C^1$]
Sea $F \colon U \subseteq \R^n \to \R^m$, con $U$ abierto. Si todas las derivadas parciales de las componentes de $F$ existen en todo punto de $U$ y las funciones
\[
\frac{\partial F_i}{\partial x_j}\colon U \to \R
\]
son continuas, entonces se dice que $F$ es una \emph{funci\'on de clase $C^1$}, o bien que es \emph{continuamente diferenciable} en $U$.
\end{definition}

\begin{remark}
Si $F$ es diferenciable en $p$, entonces para todo $v\in \R^n$ existe la derivada direccional
\[
(\dd F)_p(v)
=
\lim_{t\to 0}\frac{F(p+tv)-F(p)}{t}.
\]
En particular, existen todas las derivadas parciales y $F$ es continua en $p$.
\end{remark}

\begin{remark}
La existencia de todas las derivadas direccionales en un punto no garantiza que $F$ sea diferenciable en ese punto.
\end{remark}

\begin{example}
Consid\'erese la funci\'on $F \colon \R^2 \to \R$ dada por
\[
F(x,y)=
\begin{cases}
\dfrac{x^2y}{x^4+y^2}, & (x,y)\neq (0,0),\\[0.6em]
0, & (x,y)=(0,0).
\end{cases}
\]
Este es un ejemplo cl\'asico donde la existencia de ciertas derivadas direccionales no alcanza por s\'i sola para garantizar diferenciabilidad en el origen.
\end{example}

\end{document}
